\begin{abstract}
Vision transformers (ViTs) spend most of their compute on a fixed set of image tokens, and token-reduction methods developed on ImageNet promise large savings. Whether those savings transfer to small medical datasets, where a single missed class can matter, is unclear. We study token reduction and parameter-efficient adaptation for a \deit backbone on three MedMNIST benchmarks (\pathmnist, \dermamnist, \pneumnist) under a pre-registered protocol: budgets are chosen on validation only, test configurations are hashed and locked before any test inference, and non-inferiority is judged with conservative paired bounds rather than point estimates. Across 279 validation configurations, 90 ablations and 48 locked test evaluations we find that (i) the tolerable reduction is strongly dataset-dependent, from 90\% of tokens on \pathmnist to 25\% on \dermamnist; (ii) attention-guided deletion is not reliably better than random deletion; (iii) ToMe-style merging extends tolerance far beyond deletion on two datasets; and (iv) despite small mean losses, \emph{no} policy established the predeclared 0.5\,\pp non-inferiority margin on held-out test data across three seeds, while two extra epochs of recovery training shrank mean test losses to within 0.4\,\pp. We further compare full, partial, linear and LoRA fine-tuning for both initial adaptation and post-pruning recovery, measuring accuracy against trainable parameters, training time and memory. \TODO{one-sentence fine-tuning/LoRA result once the grid finishes.} Code, locked specifications and per-image predictions are released.
\end{abstract}
